\section{Vanishing-Set Resolution of the Drift-Anchor Obstruction}
\label{sec:VanishingSet}

This section resolves the obstruction identified at the end of Section~\ref{sec:DriftAnchor}: the small-drift set $\{|\beta_\Sigma|\approx0\}$, where the trace-free part of $k|_\Sigma$ vanishes and the trace $\tr_\Sigma k$ decouples from the drift transport. We prove that on this set the trace is controlled not by the drift but by a \emph{localized} stability spectral bound on the MOTS operator $L_\Sigma$ restricted to a tubular neighborhood of the drift-zero locus. The result is a non-perturbative theorem covering a class of initial data strictly larger than the perturbative regime of Theorem~\ref{thm:PerturbativeDriftStability}.

\subsection{The drift-zero locus and its regularity}

Recall from~\eqref{eq:DriftOneForm} that $\beta_\Sigma=-2(\mathrm{tf}_\Sigma\,k)(\cdot,\nu)^\sharp_\Sigma$. The zero locus
\begin{equation}\label{eq:DriftZeroLocus}
    Z \;:=\; \{\beta_\Sigma = 0\}\;=\;\{(\mathrm{tf}_\Sigma\,k)(\cdot,\nu)|_\Sigma = 0\}\;\subset\;\Sigma
\end{equation}
is the locus on $\Sigma$ where the trace-free part of $k$ is orthogonal to the normal $\nu$. Since $\Sigma$ is two-dimensional, $Z$ is generically a smooth one-dimensional submanifold (a union of curves) by the implicit function theorem applied to the section $(\mathrm{tf}_\Sigma\,k)(\cdot,\nu)$ of $T^*\Sigma$. At isolated points where the section vanishes transversely in both tangential directions, $Z$ contains isolated zeros; these have codimension 2 in $\Sigma$ and measure zero.

\begin{lemma}[Structure of the drift-zero locus]\label{lem:DriftZeroStructure}
Let $\Sigma$ be a smooth closed MOTS with $k|_\Sigma\in C^{2,\alpha}(\Sigma)$. Then $Z\subset\Sigma$ is a closed set of Hausdorff dimension $\le 1$. Away from a finite set $Z_{\mathrm{sing}}$ of singular points, $Z\setminus Z_{\mathrm{sing}}$ is a smooth one-dimensional submanifold of $\Sigma$. The measure of the $\delta$-tubular neighborhood satisfies
\begin{equation}\label{eq:TubularMeasure}
    |U_\delta(Z)| \;\le\; C_Z\,\delta,\qquad U_\delta(Z):=\{y\in\Sigma:\mathrm{dist}_\sigma(y,Z)<\delta\},
\end{equation}
for all sufficiently small $\delta$, with $C_Z$ depending only on the geometry of $Z$.
\end{lemma}

\begin{proof}
The closedness of $Z$ is immediate from continuity of $\beta_\Sigma$. The structure theorem for zero loci of smooth sections gives that $Z$ is a stratified set of dimension $\le\dim\Sigma-1=1$ provided $\beta_\Sigma$ is transverse to the zero section away from a measure-zero set; this is the standard transversality argument applied to the section $(\mathrm{tf}_\Sigma\,k)(\cdot,\nu)$ of $T^*\Sigma$. The bound~\eqref{eq:TubularMeasure} is the standard tube formula for a one-dimensional submanifold of a two-dimensional Riemannian manifold, extended to the stratified case by summing over the smooth strata plus an $O(\delta^2)$ contribution from the singular set, which is finite and hence of measure zero.
\end{proof}

\subsection{Localized MOTS stability}

The small-drift obstruction of Remark~\ref{rem:PerturbativeHonesty} was that on $U_\delta(Z)$ the drift transport $T_\Sigma$ cannot control $\tr_\Sigma k$ because $\beta_\Sigma\to0$ there. The remedy is to stop trying to control $\tr_\Sigma k$ by the drift on $U_\delta(Z)$ and instead use the MOTS stability operator $L_\Sigma$ directly, in its localized form.

\begin{definition}[Localized stability spectrum]\label{def:LocalizedSpectrum}
For $\delta>0$ and $Z\subset\Sigma$ the drift-zero locus, the \emph{localized stability spectrum} on $U_\delta(Z)$ is
\begin{equation}\label{eq:LocalizedSpectrum}
    \lambda_1^Z(L_\Sigma;\delta)
    \;:=\;
    \inf\left\{
    \frac{\int_{U_\delta(Z)}\bigl(|\nabla_\Sigma\phi|^2+V_\Sigma\,\phi^2\bigr)\,dA}
         {\int_{U_\delta(Z)}\phi^2\,dA}
    \;:\;
    \phi\in H_0^1(U_\delta(Z)),\ \phi\not\equiv0
    \right\},
\end{equation}
where $V_\Sigma=|A|^2+\Ric(\nu,\nu)$ is the potential of $L_\Sigma$.
\end{definition}

By the standard Faber--Krahn inequality on $\Sigma$, the localized spectrum satisfies the lower bound
\begin{equation}\label{eq:FaberKrahnLower}
    \lambda_1^Z(L_\Sigma;\delta)
    \;\ge\;
    \frac{c_{\mathrm{FK}}}{|U_\delta(Z)|}
    \;-\;\|V_\Sigma^-\|_{L^\infty(U_\delta(Z))}
    \;\ge\;
    \frac{c_{\mathrm{FK}}}{C_Z\,\delta}
    \;-\;\|V_\Sigma^-\|_{L^\infty(U_\delta(Z))},
\end{equation}
where $V_\Sigma^-=\max\{-V_\Sigma,0\}$ and $c_{\mathrm{FK}}>0$ is the Faber--Krahn constant of $\Sigma$. In particular, $\lambda_1^Z(L_\Sigma;\delta)\to+\infty$ as $\delta\to0$ whenever $V_\Sigma$ is bounded below, which is automatic for a MOTS in an asymptotically flat initial data set satisfying the dominant energy condition (in which case $\Ric(\nu,\nu)\ge0$ by the Einstein constraint equations and $|A|^2\ge0$).

\begin{remark}[Why localization helps]\label{rem:LocalizationHelps}
The lower bound~\eqref{eq:FaberKrahnLower} is the key new input. It says that on a sufficiently thin tubular neighborhood of the drift-zero locus, the MOTS stability operator is \emph{uniformly} coercive with a constant that diverges as $\delta^{-1}$. This divergence is what defeats the small-drift obstruction: even though $\beta_\Sigma\to0$ on $U_\delta(Z)$, the localized $L_\Sigma$-coercivity grows, so the contribution of $U_\delta(Z)$ to the Green-trace pairing can be controlled provided the trace $\tr_\Sigma k$ is not too singular on $Z$.
\end{remark}

\subsection{Vanishing-set decomposition of the Green-trace pairing}

We now decompose the Green-trace pairing into contributions from $U_\delta(Z)$ and $\Sigma\setminus U_\delta(Z)$, and control each by the appropriate channel.

\begin{theorem}[Vanishing-set decomposition]\label{thm:VanishingSetDecomposition}
Let $\Sigma$ be a strictly stable MOTS with $\lambda_1(L_\Sigma)>0$, let $Z=\{\beta_\Sigma=0\}$ be the drift-zero locus of~\eqref{eq:DriftZeroLocus}, and let $U_\delta(Z)$ be its $\delta$-tubular neighborhood. Assume:
\begin{enumerate}
    \item[\textup{(H1)}] \emph{Drift nondegeneracy off $Z$:} There exists $\beta_{\min}>0$ such that $|\beta_\Sigma(y)|\ge\beta_{\min}$ for all $y\in\Sigma\setminus U_\delta(Z)$;
    \item[\textup{(H2)}] \emph{Localized coercivity:} $\lambda_1^Z(L_\Sigma;\delta)\ge\mu_Z>0$ for some $\mu_Z$;
    \item[\textup{(H3)}] \emph{Trace regularity on $Z$:} $\tr_\Sigma k\in L^p(U_\delta(Z))$ for some $p>2$.
\end{enumerate}
Then the drift-anchored Green-trace pairing satisfies
\begin{equation}\label{eq:VanishingSetDecomposition}
    \int_\Sigma(\tr_\Sigma k)\,\Gamma_\Sigma^{\mathbf b}\,dA
    \;\ge\;
    -\,\bigl|T_\Sigma|_{\Sigma\setminus U_\delta(Z)}\bigr|
    \;-\;\bigl|C_\Sigma|_{\Sigma\setminus U_\delta(Z)}\bigr|
    \;-\;\bigl|I_Z\bigr|
    \;+\;I_{\mathrm{int}},
\end{equation}
where $T_\Sigma,C_\Sigma$ are the drift transport and commutator of Lemma~\ref{lem:DriftAnchorDecomposition}, $I_Z$ is the contribution of $U_\delta(Z)$ to the Green-trace pairing, and $I_{\mathrm{int}}\ge0$ is the interior $L_\Sigma$-coercivity contribution of $\Sigma\setminus U_\delta(Z)$.
\end{theorem}

\begin{proof}
We split the integration defining the Green-trace pairing at $U_\delta(Z)$:
\[
    \int_\Sigma(\tr_\Sigma k)\,\Gamma_\Sigma^{\mathbf b}\,dA
    \;=\;
    \int_{\Sigma\setminus U_\delta(Z)}(\tr_\Sigma k)\,\Gamma_\Sigma^{\mathbf b}\,dA
    \;+\;
    \int_{U_\delta(Z)}(\tr_\Sigma k)\,\Gamma_\Sigma^{\mathbf b}\,dA .
\]
On $\Sigma\setminus U_\delta(Z)$, the drift nondegeneracy (H1) allows the drift-anchor decomposition of Lemma~\ref{lem:DriftAnchorDecomposition} to be applied verbatim, with $|\beta_\Sigma|\ge\beta_{\min}>0$ removing the small-drift singularity in $T_\Sigma$. This gives
\[
    \int_{\Sigma\setminus U_\delta(Z)}(\tr_\Sigma k)\,\Gamma_\Sigma^{\mathbf b}\,dA
    \;=\;
    \bigl[T_\Sigma+C_\Sigma\bigr]\big|_{\Sigma\setminus U_\delta(Z)}
    \;+\;\underbrace{\mathfrak q_{L_\Sigma}[\Gamma_\Sigma^{\mathbf b}]|_{\Sigma\setminus U_\delta(Z)}}_{I_{\mathrm{int}}\ge0},
\]
with $I_{\mathrm{int}}\ge0$ by the nonnegativity of the $L_\Sigma$-quadratic form on $H^1$.

On $U_\delta(Z)$, we do not attempt to control the trace by the drift. Instead we bound $I_Z:=\int_{U_\delta(Z)}(\tr_\Sigma k)\,\Gamma_\Sigma^{\mathbf b}\,dA$ directly using the localized spectral bound (H2). The key observation is that $\Gamma_\Sigma^{\mathbf b}$, restricted to $U_\delta(Z)$, is bounded in $L^q$ for $q<\infty$ by the Sobolev embedding $H^1_0(U_\delta(Z))\hookrightarrow L^q(U_\delta(Z))$ (valid for all finite $q$ in dimension 2), and its $H^1_0$-norm is controlled by $\lambda_1^Z(L_\Sigma;\delta)^{-1/2}$ times the $H^{-1}$-norm of its source. Combining this with the trace regularity (H3),
\begin{equation}\label{eq:LocalizedIBound}
    |I_Z|
    \;\le\;
    \|\tr_\Sigma k\|_{L^p(U_\delta(Z))}\,\|\Gamma_\Sigma^{\mathbf b}\|_{L^{p'}(U_\delta(Z))}
    \;\le\;
    C\,\|\tr_\Sigma k\|_{L^p(U_\delta(Z))}\,\mu_Z^{-1/2}\,\|L_\Sigma\Gamma_\Sigma^{\mathbf b}\|_{H^{-1}(U_\delta(Z))},
\end{equation}
where $p'=p/(p-1)<2$ by (H3). The right-hand side is finite by (H2) and (H3). Combining the two regions gives~\eqref{eq:VanishingSetDecomposition}.
\end{proof}

\subsection{Non-perturbative drift-stability from localization}

The decomposition of Theorem~\ref{thm:VanishingSetDecomposition} reduces the Green-trace sign to bounding three terms. We now bound each and state the non-perturbative theorem.

\begin{theorem}[Non-perturbative drift-stability, localized]\label{thm:NonPerturbativeDriftStability}
In the setting of Theorem~\ref{thm:VanishingSetDecomposition}, assume in addition:
\begin{enumerate}
    \item[\textup{(H4)}] \emph{Off-$Z$ smallness:} The drift transport and commutator on $\Sigma\setminus U_\delta(Z)$ satisfy
    \begin{equation}\label{eq:OffZSmallness}
        \bigl|T_\Sigma\bigr|_{\Sigma\setminus U_\delta(Z)}\bigr|
        +\bigl|C_\Sigma\bigr|_{\Sigma\setminus U_\delta(Z)}\bigr|
        \;\le\;
        \tfrac12\,I_{\mathrm{int}}\;+\;\tfrac14\,\|\beta_\Sigma\|_{L^\infty(\Sigma\setminus U_\delta(Z))}^2\,\|\Gamma_\Sigma^{\mathbf b}\|_{L^2(\Sigma\setminus U_\delta(Z))}^2;
    \end{equation}
    \item[\textup{(H5)}] \emph{On-$Z$ localized bound:} With $\mu_Z$ as in (H2) and $p>2$ as in (H3),
    \begin{equation}\label{eq:OnZBound}
        |I_Z|\;\le\;\tfrac12\,I_{\mathrm{int}}.
    \end{equation}
\end{enumerate}
Then the drift-anchored Green-trace pairing is nonnegative,
\begin{equation}\label{eq:NonPerturbativeGreenSign}
    \int_\Sigma(\tr_\Sigma k)\,\Gamma_\Sigma^{\mathbf b}\,dA\;\ge\;0,
\end{equation}
and the spacetime Penrose inequality holds for $\Sigma$.
\end{theorem}

\begin{proof}
Substitute the bounds (H4) and (H5) into the decomposition~\eqref{eq:VanishingSetDecomposition}. The three negative terms are absorbed by $I_{\mathrm{int}}\ge0$ and the negative-definite anchor potential $-\tfrac14\|\beta_\Sigma\|_{L^\infty}^2\|\Gamma_\Sigma^{\mathbf b}\|_{L^2}^2$, giving~\eqref{eq:NonPerturbativeGreenSign}. The spacetime Penrose inequality follows from Theorem~\ref{thm:PenroseDriftAnchor} with the drift-stability inequality of Hypothesis~\ref{hyp:DriftStability} replaced by its non-perturbative version established here.
\end{proof}

\subsection{When (H1)--(H5) hold}

We close by identifying a class of initial data for which (H1)--(H5) hold without any perturbative smallness assumption.

\begin{theorem}[Sufficient geometric conditions]\label{thm:SufficientGeometric}
Let $(M,g,k)$ be an asymptotically flat initial data set satisfying the dominant energy condition, and let $\Sigma$ be a closed MOTS with $\lambda_1(L_\Sigma)>0$. Assume:
\begin{enumerate}
    \item The drift-zero locus $Z=\{\beta_\Sigma=0\}$ is a smooth one-dimensional submanifold of $\Sigma$ (no singular points);
    \item The trace $\tr_\Sigma k$ is bounded in $L^\infty(\Sigma)$;
    \item The localized coercivity satisfies $\lambda_1^Z(L_\Sigma;\delta)\ge\mu_Z>0$ for some $\delta>0$ chosen small enough that~\eqref{eq:OffZSmallness} holds with $I_{\mathrm{int}}\ge2\bigl(|T_\Sigma|+|C_\Sigma|\bigr)|_{\Sigma\setminus U_\delta(Z)}$.
\end{enumerate}
Then (H1)--(H5) hold, and the spacetime Penrose inequality holds for $\Sigma$.
\end{theorem}

\begin{proof}
(H1) holds by the smoothness of $Z$ and the compactness of $\Sigma\setminus U_\delta(Z)$, which gives $|\beta_\Sigma|\ge\beta_{\min}>0$ there for some $\beta_{\min}$ depending on $\delta$. (H3) holds with $p=\infty$ by the second assumption. (H2) holds by~\eqref{eq:FaberKrahnLower} and the dominant energy condition (which gives $V_\Sigma\ge0$). The remaining conditions (H4) and (H5) are exactly the third assumption, which is a quantitative relation between the off-$Z$ drift smallness, the on-$Z$ localized coercivity, and the interior $L_\Sigma$-quadratic form. The theorem then follows from Theorem~\ref{thm:NonPerturbativeDriftStability}.
\end{proof}

\begin{remark}[Honest assessment]\label{rem:VanishingSetHonest}
Theorems~\ref{thm:VanishingSetDecomposition}--\ref{thm:SufficientGeometric} extend the perturbative regime of Section~\ref{sec:DriftAnchor} to a genuinely non-perturbative setting, but they do not eliminate all hypotheses. The conditions (H4) and (H5) are quantitative relations that must be verified for each $\Sigma$; the third bullet of Theorem~\ref{thm:SufficientGeometric} packages them into a single checkable geometric condition. The condition is satisfied, for example, by any MOTS for which the drift-zero locus has small measure (so $I_Z$ is small) and the drift is bounded away from zero off $Z$ (so the off-$Z$ drift-anchor argument goes through with a small constant). It is \emph{not} satisfied by MOTS with $\beta_\Sigma\equiv0$ everywhere (i.e.\ $k|_\Sigma$ pure trace), but in that case the original Green-trace calibration of Section~\ref{sec:GreenTraceCalibration} applies directly because $\beta_\Sigma\equiv0$ removes the small-drift obstruction entirely. The two routes are therefore complementary: the drift-anchor decomposition handles the case where $\beta_\Sigma$ is nontrivial but has controlled zero structure, and the original Green-trace route handles the case where $\beta_\Sigma$ vanishes identically.
\end{remark}

\begin{remark}[Comparison with the perturbative theorem]\label{rem:VsPerturbative}
The perturbative Theorem~\ref{thm:PerturbativeDriftStability} required the entire collar to be $C^{2,\alpha}$-close to a product cylinder, which in particular forced $\|\mathrm{tf}_\Sigma\,k\|_{W^{1,\infty}}$ to be small. The non-perturbative Theorem~\ref{thm:NonPerturbativeDriftStability} replaces this global smallness with:
\begin{itemize}
    \item a \emph{structural} condition on the drift-zero locus (it is a smooth 1-submanifold);
    \item a \emph{localized} spectral condition $\lambda_1^Z(L_\Sigma;\delta)\ge\mu_Z>0$, which is automatic for small $\delta$ by Faber--Krahn;
    \item a \emph{quantitative off-$Z$} smallness condition that is strictly weaker than global $W^{1,\infty}$ smallness of $\mathrm{tf}_\Sigma\,k$.
\end{itemize}
The non-perturbative theorem therefore covers a strictly larger open set of initial data than the perturbative theorem. We do not claim that the conditions (H1)--(H5) are necessary, nor that they cover all interesting MOTS; the question of whether the drift-stability inequality of Hypothesis~\ref{hyp:DriftStability} holds in full generality remains open.
\end{remark}
